\section*{Chapitre \ref{ch:g12:early-earth} --- La Terre primitive et son atmosphère}

\begin{solution}{exo:g12:early-earth:1}
(a) le \omterm{def:g12:early-earth:primitive-atmosphere}{dégazage}~; (b) un \omterm{def:g12:early-earth:stromatolite}{stromatolithe}~; (c) une \omterm{def:g12:early-earth:bif}{formation de fer
rubané}~; (d) les \omterm{def:g12:early-earth:red-beds}{couches rouges}~; (e) la \omterm{def:g12:early-earth:goe}{Grande Oxydation}.
\end{solution}

\begin{solution}{exo:g12:early-earth:2}
Les zircons qui montrent de l'eau liquide (4.4~milliards d'années)~; les
premiers \omterm{def:g12:early-earth:stromatolite}{stromatolithes} du Pilbara (3.43)~; les \omterm{def:g12:early-earth:bif}{formations de fer rubané}
des monts Hamersley (plus de 2.5)~; la fin de la \omterm{def:g12:early-earth:goe}{Grande Oxydation} (2.06)~;
la fin de la plupart des \omterm{def:g12:early-earth:bif}{formations de fer rubané} (1.8).
\end{solution}

\begin{solution}{exo:g12:early-earth:3}
À $t = 0$, $L = L_\odot/(1 + 0.4) = L_\odot/1.4 \approx 0.714\,L_\odot$~:
environ \qty{29}{\percent} de moins, près d'un tiers de luminosité en moins
par rapport à aujourd'hui.
\end{solution}

\begin{solution}{exo:g12:early-earth:4}
Que l'air ne contenait alors presque pas d'oxygène. La pyrite est détruite
par l'oxygène~; ses grains n'ont pu être roulés par les rivières et déposés
intacts que si l'air et l'eau de surface en étaient dépourvus.
\end{solution}

\begin{solution}{exo:g12:early-earth:5}
La vapeur d'eau, qui s'est condensée en océan, et le dioxyde de carbone,
qui est resté dans l'air en quantités bien plus grandes qu'aujourd'hui et a
réchauffé la Terre par son effet de serre.
\end{solution}

\begin{solution}{exo:g12:early-earth:6}
Il y a 3~milliards d'années~: $1 - t/t_\odot = 3/4.567 \approx 0.657$, $L \approx
1/1.263 \approx 0.792\,L_\odot$, $T \approx 255 \times 0.792^{1/4} \approx
\qty{241}{K}$. Il y a 2.45~milliards d'années~: $1 - t/t_\odot \approx 0.536$, $L
\approx 1/1.215 \approx 0.823\,L_\odot$, $T \approx 255 \times 0.823^{1/4}
\approx \qty{243}{K}$.
\end{solution}

\begin{solution}{exo:g12:early-earth:7}
\qty{1}{kg} de fer représente $1000/55.85 \approx 17.9$ moles~; une
molécule d'oxygène prend quatre ions de fer, il faut donc
$17.9/4 \approx 4.48$ moles de
\ce{O2}, soit $4.48 \times 32.00 \approx \qty{143}{g}$ d'oxygène.
\end{solution}

\begin{solution}{exo:g12:early-earth:8}
(1) $2.46 - 2.06 = 0.40$~milliard d'années, soit 400 millions d'années.
(2) D'il y a 4.54 à il y a 2.46~milliards d'années~:
$2.08/4.54 \approx \qty{46}{\percent}$ de l'histoire de la Terre.
(3) Environ $3.43 - 2.46 \approx 1$~milliard d'années~: l'oxygène produit
par les microbes a longtemps été consommé avant de pouvoir rester dans
l'air.
\end{solution}

\begin{solution}{exo:g12:early-earth:9}
La \omterm{def:g12:early-earth:bif}{formation de fer rubané} révèle une mer profonde sans oxygène, qui
pouvait transporter du fer dissous, sous une mer peu profonde et ensoleillée
contenant assez d'oxygène pour le précipiter. La couche rouge, déposée sur
les continents, révèle de l'oxygène dans l'air et dans l'eau de surface, en
quantité suffisante pour rouiller le fer de ses grains.
\end{solution}

\begin{solution}{exo:g12:early-earth:10}
À Hamelin Pool, à Shark Bay, en Australie-Occidentale, dont l'eau est deux
fois plus salée que l'eau de mer ordinaire. Un tapis de microbes, surtout
des cyanobactéries, piège et lie des grains de sédiment et fait croître du
carbonate~; un nouveau tapis pousse par-dessus, et ainsi de suite, couche
après couche.
\end{solution}

\begin{solution}{exo:g12:early-earth:11}
(1) $21 \times 10^{-5} = \qty{0.00021}{\percent}$. (2) $21/10 =
\qty{2.1}{\percent}$.
\end{solution}

\begin{solution}{exo:g12:early-earth:12}
La \omterm{def:g10:habitable-planet:equilibrium-temperature}{température d'équilibre} est inchangée si $L(1 - A')$ est égal à $1 - A$~:
$1 - A' = 0.706/0.741 \approx 0.953$, donc $A' \approx 0.047$. La Terre
aurait dû réfléchir moins de \qty{5}{\percent} de la lumière solaire,
presque comme un corps noir, ce dont une planète océanique avec des nuages
et un peu de glace est incapable. L'autre issue est un effet de serre plus
fort.
\end{solution}

\begin{solution}{exo:g12:early-earth:13}
(1) $273 - 237 = \qty{36}{K}$ pour atteindre le gel~; $288 - 237 = \qty{51}{K}$ pour
atteindre la moyenne actuelle. (2) L'effet de serre actuel ajoute $288 - 255 =
\qty{33}{K}$~: la jeune Terre en exigeait un plus fort. (3) \qty{0.1}{bar}
représente le dixième de la pression de toute l'atmosphère actuelle, pour
le seul dioxyde de carbone.
\end{solution}

\begin{solution}{exo:g12:early-earth:14}
(1) Avant 2450 millions d'années, la signature existe~: l'air contenait
donc moins de $10^{-5}$ de l'oxygène actuel~; après 2090 millions d'années,
elle est absente~: l'air en contenait davantage. (2) Oui~: la \omterm{def:g12:early-earth:goe}{Grande
Oxydation} s'étend d'il y a environ 2460 à 2060 millions d'années, le même
intervalle.
\end{solution}

\begin{solution}{exo:g12:early-earth:15}
(1) Il y a environ $4567 - 60 = 4507$ millions d'années. (2) Au plus $4507 - 4404 =
103$ millions d'années après l'impact~: la Terre avait une croûte et de
l'eau moins d'une centaine de millions d'années après l'impact.
\end{solution}

\begin{solution}{pb:g12:early-earth:1}
\textbf{1.} $4.567 - 4.0 = 0.567$~milliard d'années.

\textbf{2.} $1 - 0.567/4.567 \approx 1 - 0.124 = 0.876$.

\textbf{3.} $L = 1/(1 + 0.4 \times 0.876) = 1/1.350 \approx 0.741$ de la
luminosité actuelle.

\textbf{4.} $0.741 \times 1361 \approx \qty{1008}{W/m^2}$.

\textbf{5.} $1 - 0.741 \approx 0.26$~: \qty{26}{\percent} de luminosité en
moins.

\textbf{6.} $T_{\mathrm{eq}} = \left(S(1-A)/(4\sigma)\right)^{1/4}$.

\textbf{7.} À $A$ et $\sigma$ fixés, $T_{\mathrm{eq}}$ est le produit d'une
constante par $S^{1/4}$, donc $T'/T = (S'/S)^{1/4}$.

\textbf{8.} $0.741^{1/4} = \sqrt{\sqrt{0.741}} \approx \sqrt{0.861} \approx
0.928$.

\textbf{9.} $255 \times 0.928 \approx \qty{237}{K}$.

\textbf{10.} $237 - 273 \approx \qty{-36}{\celsius}$.

\textbf{11.} Qu'il y avait de l'eau liquide en surface, un océan, il y a
4.4~milliards d'années.

\textbf{12.} En eau peu profonde, comme les \omterm{def:g12:early-earth:stromatolite}{stromatolithes} vivants
aujourd'hui.

\textbf{13.} Les \omterm{def:g12:early-earth:bif}{formations de fer rubané} des monts Hamersley, déposées sur
un fond marin.

\textbf{14.} Oui~: \qty{237}{K}, c'est bien en dessous de \qty{273}{K}~;
avec l'atmosphère actuelle, la surface moyenne aurait été gelée.

\textbf{15.} Les roches montrent de l'eau liquide sous un Soleil trop
faible pour la maintenir liquide avec l'atmosphère actuelle~: c'est le
\omterm{def:g12:early-earth:faint-young-sun}{paradoxe du jeune Soleil faible}.

\textbf{16.} $288 - 255 = \qty{33}{K}$.

\textbf{17.} Au moins $273 - 237 = \qty{36}{K}$, plus qu'aujourd'hui.

\textbf{18.} $A' = 1 - 0.706/0.741 \approx 0.047$~: une planète presque
noire, ce qui n'est pas plausible.

\textbf{19.} Le dioxyde de carbone, bien plus abondant qu'aujourd'hui~; le
méthane, produit par les premiers microbes, a pu y contribuer.

\textbf{20.} La \omterm{def:g10:habitable-planet:equilibrium-temperature}{température d'équilibre} de la Terre il y a quatre
milliards d'années était \textbf{d'environ 237 K}, quelque
\qty{36}{\celsius} sous le point de congélation.
\end{solution}
